\FigureLayoutDeclare{img/2009/jiangxi_science/q11_fig1.png}{16.0cm}{29bp 30bp 73bp 27bp}

\chapter{2009年江西卷（理）}

\section{单选题}

\begin{problem}
若复数 \(z = (x^2 - 1) + (x - 1)\i\) 为纯虚数，则实数 \(x\) 的值为（\quad）

\choices
    {\(-1\)}
    {\(0\)}
    {\(1\)}
    {\(-1\) 或 \(1\)}
\end{problem}

\begin{answer}
A
\end{answer}

\begin{solution}
纯虚数的实部为零且虚部不为零，因此必须满足 \(x^{2}-1=0\)，从而 \(x=-1\) 或 \(x=1\). 当 \(x=1\) 时，虚部 \(x-1=0\)，复数为零，不是纯虚数；当 \(x=-1\) 时虚部为 \(-2\)，符合条件. 所以选 A.
\end{solution}

\begin{problem}
函数 \( y = \frac{\ln(x + 1)}{\sqrt{-x^2 - 3x + 4}} \) 的定义域为（\quad）

\choices
    {\((-4, - 1)\)}
    {\((-4, 1)\)}
    {\((-1, 1)\)}
    {\((-1, 1]\)}
\end{problem}

\begin{answer}
C
\end{answer}

\begin{solution}
对数有意义要求 \(x+1>0\)，即 \(x>-1\). 分母中的根式有意义且不能为零，要求 \(-x^{2}-3x+4>0\)，即 \((x+4)(x-1)<0\)，所以 \(-4<x<1\). 两个条件取交集得定义域为 \((-1,1)\)，故选 C.
\end{solution}

\begin{problem}
已知全集 \(U = A \cup B\) 中有 \(m\) 个元素，\((\complement_U A) \cup (\complement_U B)\) 中有 \(n\) 个元素. 若 \(A \cap B\) 非空，则 \(A \cap B\) 的元素个数为（\quad）

\choices
    {\(mn\)}
    {\(m + n\)}
    {\(n - m\)}
    {\(m - n\)}
\end{problem}

\begin{answer}
D
\end{answer}

\begin{solution}
由德摩根律，\((\complement_U A)\cup(\complement_U B)=\complement_U(A\cap B)\). 因为全集 \(U\) 有 \(m\) 个元素，而 \(\complement_U(A\cap B)\) 有 \(n\) 个元素，所以

\[
|A\cap B|=|U|-|\complement_U(A\cap B)|=m-n
\]

故选 D.
\end{solution}

\begin{problem}
若函数 \(f(x) = (1 + \sqrt{3}\tan x)\cos x\)，\(0 \leqslant x < \frac{\pi}{2}\)，则 \( f(x) \) 的最大值为（\quad）

\choices
    {\(1\)}
    {\(2\)}
    {\( \sqrt{3} + 1 \)}
    {\( \sqrt{3} + 2 \)}
\end{problem}

\begin{answer}
B
\end{answer}

\begin{solution}
在给定区间内 \(\cos x>0\)，且 \(\tan x\) 有定义. 利用 \(\tan x\cos x=\sin x\)，得

\[
f(x)=\cos x+\sqrt3\sin x=2\cos\left(x-\frac\pi3\right)
\]

当 \(0\leqslant x<\frac\pi2\) 时，\(x-\frac\pi3\in\left[-\frac\pi3,\frac\pi6\right)\)，其中包含 \(0\)，所以余弦的最大值为 \(1\)，在 \(x=\frac\pi3\) 时取得. 因此 \(f(x)\) 的最大值为 \(2\)，故选 B.
\end{solution}

\begin{problem}
设函数 \( f(x) = g(x) + x^2 \)，曲线 \( y = g(x) \) 在点 \((1, g(1))\) 处的切线方程为 \( y = 2x + 1 \)，则曲线 \( y = f(x) \) 在点 \((1, f(1))\) 处切线的斜率为（\quad）

\choices
    {\(4\)}
    {\(-\frac{1}{4}\)}
    {\(2\)}
    {\(-\frac{1}{2}\)}
\end{problem}

\begin{answer}
A
\end{answer}

\begin{solution}
曲线 \(y=g(x)\) 在 \(x=1\) 处的切线斜率为 \(2\)，所以 \(g'(1)=2\). 由 \(f(x)=g(x)+x^2\)，得 \(f'(x)=g'(x)+2x\)，从而 \(f'(1)=2+2=4\). 因此所求斜率为 \(4\)，故选 A.
\end{solution}

\begin{problem}
过椭圆 \(\frac{x^2}{a^2} + \frac{y^2}{b^2} = 1\)（\(a > b > 0\)）的左焦点 \(F_1\) 作 \(x\) 轴的垂线交椭圆于点 \(P\)，\(F_2\) 为右焦点，若 \(\angle F_1PF_2 = 60^\circ\)，则椭圆的离心率为（\quad）

\choices
    {\(\frac{\sqrt{2}}{2}\)}
    {\(\frac{\sqrt{3}}{3}\)}
    {\(\frac{1}{2}\)}
    {\(\frac{1}{3}\)}
\end{problem}

\begin{answer}
B
\end{answer}

\begin{solution}
设椭圆的焦半距为 \(c\)，则 \(F_1=(-c,0)\)，\(F_2=(c,0)\)，且 \(c^2=a^2-b^2\). 由于过 \(F_1\) 的垂线为 \(x=-c\)，设 \(P=(-c,y_0)\). 代入椭圆方程得

\(\frac{c^2}{a^2}+\frac{y_0^2}{b^2}=1\)，\(|y_0|=\frac{b^2}{a}\).

在点 \(P\) 处，\(\overrightarrow{PF_1}=(0,-y_0)\)，\(\overrightarrow{PF_2}=(2c,-y_0)\)，所以

\[
\cos60^\circ=\frac{\overrightarrow{PF_1}\cdot\overrightarrow{PF_2}}{|PF_1||PF_2|}
=\frac{|y_0|}{\sqrt{4c^2+y_0^2}}
\]

于是 \(3y_0^2=4c^2\)，即 \(\frac{b^2}{a}=\frac{2c}{\sqrt3}\). 令 \(e=\frac ca\)，则

\(1-e^2=\frac{2e}{\sqrt3}\)，\(\sqrt3e^2+2e-\sqrt3=0\).

因为 \(0<e<1\)，故 \(e=\frac1{\sqrt3}=\frac{\sqrt3}{3}\). 所以选 B.
\end{solution}

\begin{problem}
\((1 + ax + by)^{n}\) 展开式中不含 \(x\) 的项的系数绝对值的和为 \(243\)，不含 \(y\) 的项的系数绝对值的和为 \(32\)，则 \(a\)，\(b\)，\(n\) 的值可能为（\quad）

\choices
    {\(a = 2, b = -1, n = 5\)}
    {\(a = -2\)，\(b = -1\)，\(n = 6\)}
    {\(a = -1\)，\(b = 2\)，\(n = 6\)}
    {\(a = 1\)，\(b = 2\)，\(n = 5\)}
\end{problem}

\begin{answer}
D
\end{answer}

\begin{solution}
展开式中不含 \(x\) 的各项，正好是令 \(x=0\) 后所得的 \((1+by)^n\). 因此这些项的系数绝对值之和为 \((1+|b|)^n\). 同理，不含 \(y\) 的各项的系数绝对值之和为 \((1+|a|)^n\). 逐项检验：当 \(a=2\)，\(b=-1\)，\(n=5\) 时，\((1+|b|)^n=2^5\ne243\)；当 \(a=-2\)，\(b=-1\)，\(n=6\) 时，\((1+|b|)^n=2^6\ne243\)；当 \(a=-1\)，\(b=2\)，\(n=6\) 时，\((1+|b|)^n=3^6\ne243\)；当 \(a=1\)，\(b=2\)，\(n=5\) 时，\((1+|b|)^n=3^5=243\)，且 \((1+|a|)^n=2^5=32\)，符合题设. 故选 D.
\end{solution}

\begin{problem}
数列 \(\{a_{n}\}\) 的通项 \(a_{n} = n^{2}\left(\cos^{2}\frac{n\pi}{3} - \sin^{2}\frac{n\pi}{3}\right)\)，其前 \(n\) 项和为 \(S_{n}\)，则 \(S_{30}\) 为（\quad）

\choices
    {\(470\)}
    {\(490\)}
    {\(495\)}
    {\(510\)}
\end{problem}

\begin{answer}
A
\end{answer}

\begin{solution}
由 \(\cos^2t-\sin^2t=\cos2t\)，得 \(a_n=n^2\cos\frac{2n\pi}{3}\). 当 \(n\) 是 \(3\) 的倍数时，\(\cos\frac{2n\pi}{3}=1\)；否则该余弦值为 \(-\frac12\). 因而

\[
S_{30}=\sum_{k=1}^{10}(3k)^2-\frac12\sum_{\substack{1\leqslant n\leqslant30\\3\nmid n}}n^2
\]

又

\(\sum_{n=1}^{30}n^2=\frac{30\times31\times61}{6}=9455\)，\(\sum_{k=1}^{10}(3k)^2=9\times385=3465\).

所以非 \(3\) 的倍数的平方和为 \(9455-3465=5990\)，从而 \(S_{30}=3465-\frac12\times5990=470\). 故选 A.
\end{solution}

\begin{problem}
如图，正四面体 \(ABCD\) 的顶点 \(A\)，\(B\)，\(C\) 分别在两两垂直的三条射线 \(Ox\)，\(Oy\)，\(Oz\) 上，则在下列命题中，错误的为（\quad）

\bitmapfigure[max width=0.42\linewidth]{img/2009/jiangxi_science/q9_fig1.png}

\choices
    {\(O - ABC\) 是正三棱锥}
    {直线 \(OB \parallel\) 平面 \(ACD\)}
    {直线 \(AD\) 与 \(OB\) 所成的角是 \(45^{\circ}\)}
    {二面角 \(D - OB - A\) 为 \(45^{\circ}\)}
\end{problem}

\begin{answer}
B
\end{answer}

\begin{solution}
设 \(OA=OB=OC=r\). 这是因为 \(AB=AC=BC\)，而三条射线两两垂直，例如 \(AB^2=OA^2+OB^2\)，其余两式同理. 按图取 \(A=(r,0,0)\)，\(B=(0,r,0)\)，\(C=(0,0,r)\)，\(D=(r,r,r)\).

三角形 \(ABC\) 的三边都为 \(\sqrt2r\)，且点 \(O\) 在平面 \(ABC\) 上的射影为该正三角形的中心，所以 \(O-ABC\) 是正三棱锥. 平面 \(ACD\) 内的两个方向向量可取 \(\overrightarrow{AC}=(-r,0,r)\) 和 \(\overrightarrow{AD}=(0,r,r)\)，其法向量可取 \(\bs{n}=(1,-1,1)\). 因为 \(\overrightarrow{OB}=(0,r,0)\) 与 \(\bs{n}\) 的内积不为零，所以 \(OB\) 不平行于平面 \(ACD\)，命题 B 错误.

又 \(\overrightarrow{AD}=(0,r,r)\)，故

\[
\cos\angle(AD,OB)=\frac{r^2}{\sqrt2r\times r}=\frac1{\sqrt2}
\]

所以直线 \(AD\) 与 \(OB\) 所成的角为 \(45^\circ\). 平面 \(AOB\) 的方程为 \(z=0\)，平面 \(DOB\) 的方程为 \(x-z=0\)，两平面的法向量夹角为 \(45^\circ\)，故二面角 \(D-OB-A\) 也为 \(45^\circ\). 因此错误的命题为 B.
\end{solution}

\begin{problem}
为了庆祝六一儿童节，某食品厂制作了 \(3\text{ 种}\) 不同的精美卡片，每袋食品随机装入一张卡片，集齐 \(3\text{ 种}\) 卡片可获奖，现购买该种食品 \(5\text{ 袋}\)，能获奖的概率为（\quad）

\choices
    {\( \frac{31}{81} \)}
    {\( \frac{33}{81} \)}
    {\( \frac{48}{81} \)}
    {\( \frac{50}{81} \)}
\end{problem}

\begin{answer}
D
\end{answer}

\begin{solution}
5 袋食品的卡片共有 \(3^5\) 种等可能结果. 设事件 \(E\) 为三种卡片均至少出现一次. 用容斥原理，缺少一种卡片的结果有 \(\mathrm C_3^1 2^5\) 种，缺少两种卡片的结果有 \(\mathrm C_3^2 1^5\) 种，所以

\[
P(E)=\frac{3^5-\mathrm C_3^1 2^5+\mathrm C_3^2 1^5}{3^5}
=\frac{243-96+3}{243}=\frac{50}{81}
\]

故选 D.
\end{solution}

\begin{problem}
一个平面封闭区域内任意两点距离的最大值称为该区域的“直径”，封闭区域边界曲线的长度与区域直径之比称为区域的“周率”，下面四个平面区域（阴影部分）的周率从左到右依次记为 \(\tau_{1}\)，\(\tau_{2}\)，\(\tau_{3}\)，\(\tau_{4}\)，则下列关系正确的为（\quad）

\bitmapfigure[max width=0.82\linewidth]{img/2009/jiangxi_science/q11_fig1.png}

\choices
    {\(\tau_{1} > \tau_{4} > \tau_{3}\)}
    {\(\tau_{3} > \tau_{1} > \tau_{2}\)}
    {\(\tau_{4} > \tau_{2} > \tau_{3}\)}
    {\(\tau_{3} > \tau_{4} > \tau_{2}\)}
\end{problem}

\begin{answer}
C
\end{answer}

\begin{solution}
记每个小正方形的边长为 \(l\). 第一图由五个小正方形组成，按图逐段相加可得边界长度为 \(12l\)；其左下角与右上角两点的横、纵坐标差均为 \(3l\)，且这两点是区域内最远的两点，所以直径为 \(3\sqrt{2}l\). 因而

\[
\tau_1=\frac{12l}{3\sqrt{2}l}=2\sqrt{2}
\]

第二图的边界由两个半径相同的半圆弧组成，弧长之和为相应圆的周长，区域直径为该圆的直径，因此 \(\tau_2=\pi\). 第三图由两个边长为 \(l\) 的正三角形沿一个顶点拼成，边界长度为 \(6l\)，最远两点间的距离为 \(2l\)，所以 \(\tau_3=3\).

第四图可看作在边长为 \(l\) 的正六边形每条边外各接一个边长为 \(l\) 的正三角形. 星形边界有 \(12\) 条长为 \(l\) 的边，两个相对外尖点的距离为 \(2\sqrt3l\)，故

\[
\tau_4=\frac{12l}{2\sqrt3l}=2\sqrt3
\]

因为 \(2\sqrt3>\pi>3>2\sqrt2\)，所以

\[
\tau_4>\tau_2>\tau_3>\tau_1
\]

因此只有选项 C 正确.
\end{solution}

\begin{problem}
设函数 \( f(x) = \sqrt{ax^2 + bx + c}\) （\(a < 0\)） 的定义域为 \(D\)，若所有点 \( (s, f(t))\) （\(s\)，\(t \in D\)） 构成一个正方形区域，则 \( a \) 的值为（\quad）

\choices
    {\(-2\)}
    {\(-4\)}
    {\(-8\)}
    {不能确定}
\end{problem}

\begin{answer}
B
\end{answer}

\begin{solution}
令 \(q(x)=ax^2+bx+c\). 因为 \(a<0\)，且定义域对应 \(q(x)\geqslant0\)，设 \(q\) 的最大值为 \(H>0\)，则可写成 \(q(x)=-\alpha(x-h)^2+H\)，其中 \(\alpha=-a>0\). 因此

\(D=\left[h-\sqrt{\frac H\alpha},h+\sqrt{\frac H\alpha}\right]\)，\(|D|=2\sqrt{\frac H\alpha}\).

函数值域为 \([0,\sqrt H]\). 所有点 \( (s,f(t)) \) 构成的区域是 \(D\times[0,\sqrt H]\)，其横向边长为 \(2\sqrt{H/\alpha}\)，纵向边长为 \(\sqrt H\). 该区域为正方形，所以

\(2\sqrt{\frac H\alpha}=\sqrt H\)，\(\frac4\alpha=1\)，\(\alpha=4\).

故 \(a=-\alpha=-4\)，选 B.
\end{solution}

\section{填空题}

\begin{problem}
已知向量 \(\bs{a} = (3, 1)\)，\(\bs{b} = (1, 3)\)，\(\bs{c} = (k, 7)\). 若 \((\bs{a} - \bs{c}) \parallel \bs{b}\)，则 \(k =\fillinblank{}\).
\end{problem}

\begin{answer}
\(5\)
\end{answer}

\begin{solution}
由 \(\bs{a}-\bs{c}=(3-k,-6)\)，而 \(\bs{b}=(1,3)\)，两向量平行当且仅当对应坐标的行列式为零，即

\[
\begin{vmatrix}3-k&-6\\1&3\end{vmatrix}=3(3-k)+6=0
\]

解得 \(k=5\).
\end{solution}

\begin{problem}
正三棱柱 \(ABC - A_{1}B_{1}C_{1}\) 内接于半径为 \(2\) 的球，若 \(A\)，\(B\) 两点的球面距离为 \(\pi\)，则正三棱柱的体积为\fillinblank{}.
\end{problem}

\begin{answer}
\(8\)
\end{answer}

\begin{solution}
设球心为 \(O\)，上下底面中心的连线为轴，设正三棱柱底面外接圆半径为 \(r\)，高为 \(h\). 设下底面中心为 \(O_1\)，则 \(OO_1\perp\) 底面，且 \(OA=OB=2\). 球面距离 \(AB=\pi\) 意味着圆心角 \(\angle AOB=\frac\pi2\)，由余弦定理得

\[
AB^2=OA^2+OB^2-2OA\cdot OB\cos\frac\pi2=8
\]

底面 \(ABC\) 是正三角形，所以 \(r=\frac{AB}{\sqrt3}=2\sqrt{\frac23}\). 在直角三角形 \(OO_1A\) 中，

\(OO_1=\sqrt{4-r^2}=\sqrt{\frac43}=\frac2{\sqrt3}\)，\(h=2OO_1=\frac4{\sqrt3}\).

底面面积为 \(\frac{\sqrt3}{4}AB^2=2\sqrt3\)，因此正三棱柱的体积为

\[
V=2\sqrt3\times\frac4{\sqrt3}=8
\]
\end{solution}

\begin{problem}
若不等式 \(\sqrt{9 - x^2} \leqslant k(x + 2) - \sqrt{2}\) 的解集为区间 \([a, b]\)，且 \(b - a = 2\)，则 \(k = \)\fillinblank{}.
\end{problem}

\begin{answer}
\(\sqrt{2}\)
\end{answer}

\begin{solution}
根式要求 \(-3\leqslant x\leqslant3\)，且右端必须非负. 若 \(k=0\)，右端恒为 \(-\sqrt2\)，无解. 若 \(k<0\)，右端非负要求

\[
x\leqslant x_0=\frac{\sqrt2}{k}-2<-2
\]

若 \(x_0<-3\)，则无解；若 \(-3\leqslant x_0<-2\)，所有解都包含在 \([-3,x_0]\) 内，该区间长度小于 \(1\)，不可能等于 \(b-a=2\). 因而 \(k>0\).

此时右端 \(k(x+2)-\sqrt2\) 随 \(x\) 增大而增大. 若 \(x_0=\frac{\sqrt2}{k}-2>3\)，则无解；所以由题设解集非空可知 \(x_0\leqslant3\)，从而 \(x=3\) 满足原不等式. 因为解集是区间 \([a,b]\)，故 \(b=3\)，再由 \(b-a=2\) 得 \(a=1\).

在右端非负的范围内平方，得

\[
9-x^2\leqslant\left(k(x+2)-\sqrt2\right)^2
\]

由于 \(x=1\) 是解集的左端点，且若右端在 \(x=1\) 处为零则原不等式左边为 \(2\sqrt2>0\)，不可能成立，所以在 \(x=1\) 处两边相等：

\(\sqrt{9-1}=k(1+2)-\sqrt2\)，\(2\sqrt2=3k-\sqrt2\)

所以 \(k=\sqrt2\). 反向验证：当 \(k=\sqrt2\) 时，右端为 \(\sqrt2(x+1)\)，故必须 \(x\geqslant-1\)，平方后得

\[
9-x^2\leqslant2(x+1)^2
\Longleftrightarrow (3x+7)(x-1)\geqslant0
\]

与 \([-1,3]\) 取交集，解集恰为 \([1,3]\)，长度为 \(2\). 所以填 \(\sqrt2\).
\end{solution}

\begin{problem}
设直线系 \(M: x \cos \theta + (y - 2) \sin \theta = 1\)（\(0 \leqslant \theta \leqslant 2\pi\)），对于下列四个命题：

\begin{circlelist}
\item \(M\) 中所有直线均经过一个定点；
\item 存在定点 \(P\) 不在 \(M\) 中的任一条直线上；
\item 对于任意整数 \(n\)（\(n \geqslant 3\)），存在正 \(n\) 边形，其所有边均在 \(M\) 中的直线上；
\item \(M\) 中的直线所能围成的正三角形面积都相等.
\end{circlelist}

其中真命题的代号是\fillinblank{}.（写出所有真命题的代号）
\end{problem}

\begin{answer}
BC
\end{answer}

\begin{solution}
令圆心为 \(O_{1}=(0,2)\)，则直线系中的直线可写为
\(\overrightarrow{O_{1}X}\cdot(\cos\theta,\sin\theta)=1\)，所以 \(M\) 中的所有直线都是圆 \(x^{2}+(y-2)^{2}=1\) 的切线.

取 \(\theta=0\) 和 \(\theta=\pi\)，所得两条直线分别为 \(x=1\) 和 \(x=-1\)，它们平行，因此所有直线不可能经过同一个定点，第一个命题错误. 圆心 \(O_{1}\) 到 \(M\) 中任一直线的距离为 \(1\)，且 \(O_{1}\) 不在这些直线上，所以第二个命题正确.

对任意整数 \(n\geqslant3\)，取 \(n\) 条法向量方向依次相差 \(\frac{2\pi}{n}\) 的切线，即取 \(\theta_{j}=\theta_{0}+\frac{2j\pi}{n}\)（\(j=0\)，\(1\)，\(\ldots\)，\(n-1\)）. 这些切线围成一个以该圆为内切圆的正 \(n\) 边形，所以第三个命题正确.

第四个命题错误. 例如，取 \(\theta=0\)，\(\frac{2\pi}{3}\)，\(\frac{4\pi}{3}\)，三条直线围成的正三角形边长为 \(2\sqrt{3}\)，面积为 \(3\sqrt{3}\). 取 \(\theta=0\)，\(\frac{\pi}{3}\)，\(\frac{2\pi}{3}\)，三条直线也围成正三角形，但其边长为 \(\frac{2}{\sqrt{3}}\)，面积为 \(\frac{\sqrt{3}}{3}\). 两个面积不相等，所以真命题的代号为 BC.
\end{solution}

\section{解答题}

\begin{problem}
设函数 \( f(x) = \frac{\e^x}{x} \).

\begin{enumerate}
\item 求函数 \( f(x) \) 的单调区间；
\item 若 \(k > 0\)，求不等式 \(f'(x) + k(1 - x)f(x) > 0\) 的解集.
\end{enumerate}
\end{problem}

\begin{solution}
\begin{enumerate}
\item
函数的定义域为 \(x\neq0\)，且
\[
f'(x)=\frac{x\e^{x}-\e^{x}}{x^{2}}=\frac{\e^{x}(x-1)}{x^{2}}
\]
因为 \(\e^{x}>0\)，\(x^{2}>0\)，所以当 \(x<1\) 且 \(x\neq0\) 时 \(f'(x)<0\)，当 \(x>1\) 时 \(f'(x)>0\). 因此 \(f(x)\) 在 \((-\infty,0)\) 和 \((0,1)\) 上单调递减，在 \((1,+\infty)\) 上单调递增.

\item
又
\[
f'(x)+k(1-x)f(x)=\frac{\e^{x}(x-1)}{x^{2}}+k(1-x)\frac{\e^{x}}{x}
=\frac{\e^{x}(x-1)(1-kx)}{x^{2}}
\]
由于 \(\frac{\e^{x}}{x^{2}}>0\)，原不等式等价于 \((x-1)(1-kx)>0\)，且 \(x\neq0\). 当 \(0<k<1\) 时，\(1<\frac1k\)，解集为 \(\left(1,\frac1k\right)\)；当 \(k=1\) 时，左侧为 \(-(x-1)^{2}\)，无解；当 \(k>1\) 时，\(\frac1k<1\)，解集为 \(\left(\frac1k,1\right)\). 综上，解集为
\[
\begin{cases}
\left(1,\frac1k\right),&0<k<1,\\
\varnothing,&k=1,\\
\left(\frac1k,1\right),&k>1
\end{cases}
\]
\end{enumerate}
\end{solution}

\begin{problem}
某公司拟资助三位大学生自主创业，现聘请两位专家，独立地对每位大学生的创业方案进行评审. 假设评审结果为“支持”或“不支持”的概率都是 \(\frac{1}{2}\)，若某人获得两个“支持”，则给予 \(10\text{ 万元}\) 的创业资助. 若只获得一个“支持”，则给予 \(5\text{ 万元}\) 的资助；若未获得“支持”，则不予资助，令 \(\xi\) 表示该公司的资助总额.

\begin{enumerate}
\item 写出 \(\xi\) 的分布列；
\item 求数学期望 \(E(\xi)\).
\end{enumerate}
\end{problem}

\begin{solution}
设第 \(i\) 位大学生获得的“支持”次数为 \(X_{i}\)（\(i=1\)，\(2\)，\(3\)）. 对每位大学生，两个专家的评审结果相互独立，故 \(P(X_{i}=0)=\frac14\)，\(P(X_{i}=1)=\frac12\)，\(P(X_{i}=2)=\frac14\).

\begin{enumerate}
\item
以下表中金额的数值单位均为万元. 每获得一次“支持”对应 \(5\) 万元资助，所以公司的资助总额满足 \(\xi=5(X_{1}+X_{2}+X_{3})\). 六次评审相互独立，令 \(K=X_{1}+X_{2}+X_{3}\)，则 \(K\) 服从参数为 \(6\)，\(\frac12\) 的二项分布. 因而
\(P(\xi=5j)=P(K=j)=\mathrm C_{6}^{j}\left(\frac12\right)^{6}=\frac{\mathrm C_{6}^{j}}{64}\)，\(j=0\)，\(1\)，\(\ldots\)，\(6\)
\begin{center}
\begin{tabular}{c|ccccccc}
\(\xi\) & \(0\) & \(5\) & \(10\) & \(15\) & \(20\) & \(25\) & \(30\) \\
\hline
\(P\) & \(\dfrac{1}{64}\) & \(\dfrac{6}{64}\) & \(\dfrac{15}{64}\) & \(\dfrac{20}{64}\) & \(\dfrac{15}{64}\) & \(\dfrac{6}{64}\) & \(\dfrac{1}{64}\)
\end{tabular}
\end{center}

\item
因此
\[
E(\xi)=5E(K)=5\times6\times\frac12=15\text{ 万元}
\]
\end{enumerate}
\end{solution}

\begin{problem}
\(\triangle ABC\) 中，\(A\)，\(B\)，\(C\) 所对的边分别为 \(a\)，\(b\)，\(c\)，\(\tan C = \frac{\sin A + \sin B}{\cos A + \cos B}\)，\(\sin(B - A) = \cos C\).

\begin{enumerate}
\item 求 \(A\)，\(C\)；
\item 若 \(S_{\triangle ABC} = 3 + \sqrt{3}\)，求 \(a\)，\(c\).
\end{enumerate}
\end{problem}

\begin{solution}
\begin{enumerate}
\item
由三角形内角和得 \(A+B=\pi-C\). 由和差化积公式，有
\(\sin A+\sin B=2\cos\frac{C}{2}\cos\frac{B-A}{2}\)，且
\(\cos A+\cos B=2\sin\frac{C}{2}\cos\frac{B-A}{2}\).
因为 \(A\)，\(B\) 是三角形内角，所以 \(\cos\frac{B-A}{2}>0\)，从而
\[
\tan C=\frac{\cos(C/2)}{\sin(C/2)}=\cot\frac{C}{2}>0
\]
故 \(0<C<\frac\pi2\). 令 \(t=\tan\frac{C}{2}>0\)，则
\[
\tan C=\frac{2t}{1-t^{2}}=\frac1t
\]
所以 \(3t^{2}=1\)，即 \(t=\frac1{\sqrt3}\)，故 \(C=\frac\pi3\).

又由已知 \(\sin(B-A)=\cos C=\frac12\). 因为 \(\lvert B-A\rvert<A+B=\frac{2\pi}{3}\)，在此范围内只能有 \(B-A=\frac\pi6\). 因此
\(A=\frac{(A+B)-(B-A)}2=\frac\pi4\)，\(B=\frac{(A+B)+(B-A)}2=\frac{5\pi}{12}\)

\item
设三角形外接圆半径为 \(R\). 由正弦定理，\(a=2R\sin A=\sqrt2R\)，\(c=2R\sin C=\sqrt3R\).
又
\[
S_{\triangle ABC}=\frac12ac\sin B
=\frac12\cdot\sqrt2R\cdot\sqrt3R\cdot\frac{\sqrt6+\sqrt2}{4}
=\frac{3+\sqrt3}{4}R^{2}
\]
由题设面积为 \(3+\sqrt3\)，得 \(R^{2}=4\)，所以 \(R=2\). 因此
\(a=2\sqrt2\)，\(c=2\sqrt3\)
\end{enumerate}
\end{solution}

\begin{problem}
在四棱锥 \(P - ABCD\) 中，底面 \(ABCD\) 是矩形，\(PA \perp\) 平面 \(ABCD\). \(PA = AD = 4\)，\(AB = 2\). 以 \(AC\) 的中点 \(O\) 为球心，\(AC\) 为直径的球面交 \(PD\) 于点 \(M\)，交 \(PC\) 于点 \(N\).

\begin{enumerate}
\item 求证：平面 \(ABM \perp\) 平面 \(PCD\)；
\item 求直线 \(CD\) 与平面 \(ACM\) 所成的角的大小；
\item 求点 \(N\) 到平面 \(ACM\) 的距离.
\end{enumerate}

\bitmapfigure[max width=0.42\linewidth]{img/2009/jiangxi_science/q20_fig1.png}
\end{problem}

\begin{solution}
建立空间直角坐标系，取 \(A\) 为原点，\(AB\)，\(AD\)，\(AP\) 分别为 \(x\)，\(y\)，\(z\) 轴正方向，则
\(A=(0,0,0)\)，\(B=(2,0,0)\)，\(D=(0,4,0)\)，\(C=(2,4,0)\)，\(P=(0,0,4)\).
球心 \(O\) 为 \(AC\) 的中点，所以 \(O=(1,2,0)\)，且
\(OA^{2}=1^{2}+2^{2}=5\). 球面方程为
\[
(x-1)^{2}+(y-2)^{2}+z^{2}=5
\]

直线段 \(PD\) 上的点可表示为 \(X(t)=(0,4t,4-4t)\)（\(0\leqslant t\leqslant1\)）. 代入球面方程得
\[
1+(4t-2)^{2}+(4-4t)^{2}=5
\Longleftrightarrow 32t^{2}-48t+16=0
\Longleftrightarrow 2t^{2}-3t+1=0
\Longleftrightarrow (2t-1)(t-1)=0
\]
其中 \(t=1\) 对应点 \(D\)，所以另一交点为 \(M=X\left(\frac12\right)=(0,2,2)\).

同理，直线段 \(PC\) 上的点可表示为 \(Y(t)=(2t,4t,4-4t)\). 代入球面方程得
\[
(2t-1)^{2}+(4t-2)^{2}+(4-4t)^{2}=5
\Longleftrightarrow 36t^{2}-52t+16=0
\Longleftrightarrow 9t^{2}-13t+4=0
\Longleftrightarrow (9t-4)(t-1)=0
\]
其中 \(t=1\) 对应点 \(C\)，所以另一交点为 \(N=Y\left(\frac49\right)=\left(\frac89,\frac{16}{9},\frac{20}{9}\right)\).

\begin{enumerate}
\item
在平面 \(ABM\) 内取
\(\overrightarrow{AB}=(2,0,0)\)，\(\overrightarrow{AM}=(0,2,2)\)，其法向量可取
\[
\bs{n}_{1}=\overrightarrow{AB}\times\overrightarrow{AM}=(0,-4,4)
\]
在平面 \(PCD\) 内取
\(\overrightarrow{PC}=(2,4,-4)\)，\(\overrightarrow{PD}=(0,4,-4)\)，其法向量可取
\[
\bs{n}_{2}=\overrightarrow{PC}\times\overrightarrow{PD}=(0,8,8)
\]
由于 \(\bs{n}_{1}\cdot\bs{n}_{2}=0\)，两平面的法向量垂直，所以平面 \(ABM\perp\) 平面 \(PCD\).

\item
由 \(\overrightarrow{AC}=(2,4,0)\) 和 \(\overrightarrow{AM}=(0,2,2)\)
可得平面 \(ACM\) 的一个法向量为
\[
\bs{n}=\overrightarrow{AC}\times\overrightarrow{AM}=(8,-4,4)=4(2,-1,1)
\]
因此平面 \(ACM\) 的方程为 \(2x-y+z=0\). 直线 \(CD\) 的方向向量可取 \(\bs{v}=(1,0,0)\). 设直线 \(CD\) 与平面 \(ACM\) 所成的角为 \(\alpha\)，则
\[
\sin\alpha=\frac{|\bs{v}\cdot(2,-1,1)|}{|\bs{v}|\sqrt{2^{2}+(-1)^{2}+1^{2}}}
=\frac{2}{\sqrt6}=\frac{\sqrt6}{3}
\]
所以 \(\tan\alpha=\sqrt2\)，即 \(\alpha=\arctan\sqrt2\).

\item
由点到平面的距离公式，
\[
d\left(N,\text{平面 }ACM\right)
=\frac{\left|2\cdot\frac89-\frac{16}{9}+\frac{20}{9}\right|}{\sqrt{2^{2}+(-1)^{2}+1^{2}}}
=\frac{20}{9\sqrt6}
=\frac{10\sqrt6}{27}
\]
\end{enumerate}
\end{solution}

\begin{problem}
已知点 \(P_{1}(x_{0}, y_{0})\) 为双曲线 \(\frac{x^{2}}{8b^{2}} - \frac{y^{2}}{b^{2}} = 1\)（\(b\) 为正常数）上任一点，\(F_{2}\) 为双曲线的右焦点. 过 \(P_{1}\) 作右准线的垂线，垂足为 \(A\)，连接 \(F_{2}A\) 并延长交 \(y\) 轴于 \(P_{2}\).

\begin{enumerate}
\item 求线段 \(P_{1}P_{2}\) 的中点 \(P\) 的轨迹 \(E\) 的方程；
\item 设轨迹 \(E\) 与 \(x\) 轴交于 \(B\)，\(D\) 两点，在 \(E\) 上任取一点 \(Q(x_{1}, y_{1})\)（\(y_{1} \neq 0\)），直线 \(QB\)，\(QD\) 分别交 \(y\) 轴于 \(M\)，\(N\) 两点. 求证：以 \(MN\) 为直径的圆过两定点.
\end{enumerate}

\bitmapfigure[max width=0.42\linewidth]{img/2009/jiangxi_science/q21_fig1.png}
\end{problem}

\begin{solution}
\begin{enumerate}
\item
将双曲线写为
\(\frac{x^{2}}{(2\sqrt2b)^{2}}-\frac{y^{2}}{b^{2}}=1\). 因此其半实轴为 \(2\sqrt2b\)，焦半距为
\(c=\sqrt{(2\sqrt2b)^{2}+b^{2}}=3b\)，右准线方程为
\(x=\frac{(2\sqrt2b)^{2}}{3b}=\frac{8b}{3}\)，右焦点为 \(F_{2}=(3b,0)\).

由题意，垂足为 \(A=\left(\frac{8b}{3},y_{0}\right)\). 直线 \(F_{2}A\) 的斜率为
\[
\frac{y_{0}}{\frac{8b}{3}-3b}=-\frac{3y_{0}}{b}
\]
所以它在 \(y\) 轴上的截距为
\[
y_{P_{2}}=-\frac{3y_{0}}{b}(0-3b)=9y_{0}
\]
即 \(P_{2}=(0,9y_{0})\). 设 \(P_{1}P_{2}\) 的中点为 \(P=(x,y)\)，则
\(x=\frac{x_{0}}2\)，\(y=5y_{0}\)，即 \(x_{0}=2x\)，\(y_{0}=\frac{y}{5}\). 代入 \(P_{1}\) 所在双曲线，得
\[
\frac{(2x)^{2}}{8b^{2}}-\frac{(y/5)^{2}}{b^{2}}=1
\]
所以轨迹 \(E\) 的方程为
\[
\frac{x^{2}}{2b^{2}}-\frac{y^{2}}{25b^{2}}=1
\]

\item
由轨迹方程令 \(y=0\)，可取
\(B=(\sqrt2b,0)\)，\(D=(-\sqrt2b,0)\). 设 \(Q=(x_{1},y_{1})\)，其中 \(y_{1}\neq0\). 直线 \(QB\) 与 \(y\) 轴交于 \(M=(0,m)\)，由两点式得
\[
m=\frac{\sqrt2b\,y_{1}}{\sqrt2b-x_{1}}
\]
同理，直线 \(QD\) 与 \(y\) 轴交于 \(N=(0,n)\)，且
\[
n=\frac{\sqrt2b\,y_{1}}{\sqrt2b+x_{1}}
\]
由 \(Q\in E\)，有
\[
\frac{x_{1}^{2}}{2b^{2}}-\frac{y_{1}^{2}}{25b^{2}}=1
\]
从而 \(2b^{2}-x_{1}^{2}=-\frac{2y_{1}^{2}}{25}\). 因此
\[
m+n=\frac{4b^{2}y_{1}}{2b^{2}-x_{1}^{2}}=-\frac{50b^{2}}{y_{1}}
\]
并且
\[
mn=\frac{2b^{2}y_{1}^{2}}{2b^{2}-x_{1}^{2}}=-25b^{2}
\]
以 \(MN\) 为直径的圆方程为
\[
x^{2}+(y-m)(y-n)=0
\]
即
\[
x^{2}+y^{2}-(m+n)y+mn=0
\]
代入上式得到
\[
x^{2}+y^{2}+\frac{50b^{2}}{y_{1}}y-25b^{2}=0
\]
令 \(y=0\)，得 \(x^{2}=25b^{2}\)，所以该圆恒过两个定点
\((5b,0)\) 和 \((-5b,0)\).
\end{enumerate}
\end{solution}

\begin{problem}
各项均为正数的数列 \(\{a_{n}\}\)，\(a_{1} = a\)，\(a_{2} = b\)，且对满足 \(m + n = p + q\) 的正整数 \(m\)，\(n\)，\(p\)，\(q\) 都有 \(\frac{a_{m} + a_{n}}{(1 + a_{m})(1 + a_{n})} = \frac{a_{p} + a_{q}}{(1 + a_{p})(1 + a_{q})}\).

\begin{enumerate}
\item 当 \(a = \frac{1}{2}\)，\(b = \frac{4}{5}\) 时，求通项 \(a_{n}\)；
\item 证明：对任意 \(a\)，存在与 \(a\) 有关的常数 \(\lambda\)，使得对于每个正整数 \(n\)，都有 \(\frac{1}{\lambda} \leqslant a_n \leqslant \lambda\).
\end{enumerate}
\end{problem}

\begin{solution}
令
\[
u_{n}=\frac{a_{n}-1}{a_{n}+1}
\]
因为 \(a_{n}>0\)，所以 \(-1<u_{n}<1\)，且
\[
a_{n}=\frac{1+u_{n}}{1-u_{n}}
\]
又有 \(1+a_{n}=\frac{2}{1-u_{n}}\)，所以
\[
\frac{a_{m}+a_{n}}{(1+a_{m})(1+a_{n})}
=\frac{(1+u_{m})(1-u_{n})+(1+u_{n})(1-u_{m})}{4}
=\frac{1-u_{m}u_{n}}{2}
\]
题设条件因此等价于：当 \(m+n=p+q\) 时，\(u_{m}u_{n}=u_{p}u_{q}\).

先取 \((m,n)=(1,n+1)\) 和 \((p,q)=(2,n)\)，得到
\(u_{1}u_{n+1}=u_{2}u_{n}\)，\((n=1,2,\ldots)\).
若 \(u_{1}\neq0\)，令 \(r=\frac{u_{2}}{u_{1}}\)，则
\(u_{n+1}=ru_{n}\)，\(u_{n}=u_{1}r^{n-1}\).
\begin{enumerate}
\item
当 \(a=\frac12\)，\(b=\frac45\) 时，
\[
u_{1}=\frac{\frac12-1}{1+\frac12}=-\frac13
\]
同理，\(u_{2}=\frac{\frac45-1}{1+\frac45}=-\frac19\).
所以 \(r=\frac13\)，从而 \(u_{n}=-\frac{1}{3^{n}}\). 因此
\[
a_{n}=\frac{1+u_{n}}{1-u_{n}}=\frac{1-\frac1{3^{n}}}{1+\frac1{3^{n}}}
=\frac{3^{n}-1}{3^{n}+1}
\]

\item
若 \(u_{1}\neq0\)，由于所有 \(u_{n}\) 都满足 \(\lvert u_{n}\rvert<1\)，不可能有 \(\lvert r\rvert>1\)，否则当 \(n\) 足够大时会有 \(\lvert u_{1}r^{n-1}\rvert\geqslant1\). 所以 \(\lvert r\rvert\leqslant1\)，从而
\(\lvert u_{n}\rvert\leqslant\lvert u_{1}\rvert\).

若 \(u_{1}=0\)，则对任意正整数 \(n\)，比较 \((m,n)=(1,2n-1)\) 与 \((p,q)=(n,n)\)，有
\[
u_{1}u_{2n-1}=u_{n}^{2}
\]
因此 \(u_{n}=0\)，即 \(a_{n}=1\). 这时结论显然成立.

当 \(u_{1}\neq0\) 时，由 \(a_{n}=\frac{1+u_{n}}{1-u_{n}}\) 且该函数在 \((-1,1)\) 上单调递增，得
\[
\frac{1-\lvert u_{1}\rvert}{1+\lvert u_{1}\rvert}
\leqslant a_{n}\leqslant
\frac{1+\lvert u_{1}\rvert}{1-\lvert u_{1}\rvert}
\]
又
\[
\frac{1+\lvert u_{1}\rvert}{1-\lvert u_{1}\rvert}
=\frac{a+1+\lvert a-1\rvert}{a+1-\lvert a-1\rvert}
=\begin{cases}
a,&a\geqslant1,\\
\frac1a,&0<a<1,
\end{cases}
=\max\left\{a,\frac1a\right\}
\]
取
\(\lambda=\max\left\{a,\frac1a\right\}\)，则对每个正整数 \(n\) 都有
\[
\frac1\lambda\leqslant a_{n}\leqslant\lambda
\]
这也包含 \(a=1\) 时的情形.
\end{enumerate}
\end{solution}
